\subsection{Lodging Industry}

\paragraph{Scale and significance.}
Providing overnight accommodation ``\emph{goes back into antiquity; it is the world's oldest commercial business}''~\cite{goeldner2009} (p.~180). The AH\&LA 2007 Lodging Industry Profile cited in the textbook quantifies the scale of the U.S.~sector alone: ``\emph{47{,}135 properties, 4.4 million rooms, \$133.4 billion in sales, \$26.6 billion in pretax profit, 63.3 percent average occupancy rate, and \$61.93 revenue per available room}'' (p.~156). An average occupancy of 63.3\% means that more than one-third of all hotel rooms go unsold every night---a substantial revenue gap.

\paragraph{Key performance indicators.}
The textbook's glossary defines the three metrics that every revenue manager monitors (p.~186):
\begin{itemize}
  \item \textbf{ADR} (Average Daily Rate): ``\emph{guestroom revenue for a given period divided by total number of occupied rooms.}''
  \item \textbf{Occupancy}: ``\emph{available rooms occupied for a given period. Computed by dividing the number of rooms occupied by the number of rooms available.}''
  \item \textbf{RevPAR}: ``\emph{guestroom revenue for a given period divided by total number of available rooms. A simple way to calculate is to multiply the occupancy times the room rate.}''
\end{itemize}
These three metrics are tightly coupled: raising rates (higher ADR) may reduce occupancy, while lowering rates to fill rooms may reduce ADR. RevPAR is the composite measure that captures this trade-off and is the primary output that our optimisation model targets.

\paragraph{Customer segments.}
``\emph{Who is the typical lodging customer? According to the AH\&LA, 44 percent are business travelers and 56 percent are leisure travelers}'' (p.~156). The typical business room-night is generated by a male (65\%) aged 35--54, earning \$85{,}900 per year, travelling alone (56\%), and paying \$112 per night. The typical leisure room-night involves two adults, earning \$77{,}100, travelling by automobile (77\%), and paying \$103 per night. The difference in price sensitivity and booking behaviour between these two segments is the basis for the group-versus-transient allocation decision in our model.

\paragraph{Perishable inventory, fixed capacity, and seasonality.}
Three structural characteristics make revenue management both necessary and difficult:
\begin{enumerate}
  \item \textbf{Perishability}: an unsold room-night cannot be stored or resold. This is the same property that drove the airline industry to originate yield management~\cite{rothstein1971,talluri2004}.
  \item \textbf{Fixed capacity}: ``\emph{accommodations must be available in sufficient quantity to match the demand of the travelers who arrive at the destination}'' (p.~388), but a hotel cannot add or remove rooms on a daily basis.
  \item \textbf{Seasonality}: demand fluctuates predictably across the year while supply remains constant, creating excess demand in peak periods and excess capacity in off-peak periods.
\end{enumerate}

\paragraph{Industry response: yield management.}
The textbook records the industry's response as a clear trend: ``\emph{Other trends are the increased use of central reservation systems, emphasis on service, and the use of \textbf{yield management techniques}}'' (p.~190). The glossary defines yield management as ``\emph{the use of pricing and inventory controls, based upon historical data, to maximize profits by offering varying fares over time for the same product}'' (p.~698). Contemporary industry guidance from Amadeus~\cite{amadeus2023} confirms that revenue management now rests on four foundational principles---\emph{segmentation}, \emph{forecasting}, \emph{optimisation}, and \emph{distribution}---which map directly onto the models we develop.
The canonical statement of the goal---selling ``\emph{the right [room] to the right customers and at the right time}''---comes from the 1987 annual report of American Airlines, the company that pioneered the practice~\cite{anderson2012,smith1992}.
The economics justify the effort: hotel revenue management ``\emph{has been shown to increase revenue by 2 to 5 percent}'', and because the industry's cost structure is dominated by fixed costs, ``\emph{a large portion of this revenue increase flows directly to the bottom line}''~\cite{anderson2012}; commercial revenue management decision support systems report revenue gains of approximately 6\%~\cite{boyd2003}, and American Airlines' pioneering system 4--7\%~\cite{smith1992}.
The research literature states the same idea in the language of operations research: revenue management is the utilisation of perishable assets to maximise revenues over a short-term horizon given that costs are effectively fixed over that horizon~\cite{mcgill1999}---confirming the first two characteristics above (perishability, and the fixed-cost structure that fixed capacity implies) from the operations-research tradition; seasonality is grounded in the tourism-demand analysis of Chapter~13 of the textbook.
